% --- Archon physics formalization source begin ---
% archon:physics
% archon:covers PhyXMiniProblems/problem_phyx_mini_0709.lean
% archon:source-report reports/phyx_mini/problem_phyx_mini_0709.source.json

\chapter{Physics problem phyx\_mini\_0709}
\label{ch:PhyXMiniProblems_problem_phyx_mini_0709}

\paragraph{Problem source.}
\#\# Physical scenario

A \textbackslash{}(3.0\textbackslash{},\textbackslash{}mathrm\{m\}\textbackslash{})-long ladder leans against a frictionless wall at an angle of \textbackslash{}(60\textasciicircum{}\textbackslash{}circ\textbackslash{}). The coefficient of static friction with the ground, that prevents the ladder from slipping.

\#\# Question

What is the minimum value of \textbackslash{}(\textbackslash{}mu\_s\textbackslash{})?

\#\# Answer choices

- A: A: 0.35

- B: B: 0.23

- C: C: 0.29

- D: D: 0.41

\#\# Figure caption (auxiliary; use the image as primary evidence)

The image depicts a static equilibrium scenario involving a rod leaning against a wall. Here's a detailed description of the content:

1. **Rod and Dimensions**:
   - A rod of length \textbackslash{}( L = 3.0 \textbackslash{}, \textbackslash{}text\{m\} \textbackslash{}) is shown leaning against a vertical wall.
   - The rod forms a \textbackslash{}( 60\textasciicircum{}\textbackslash{}circ \textbackslash{}) angle with the horizontal surface.

2. **Forces**:
   - **\textbackslash{}( \textbackslash{}vec\{F\}\_G \textbackslash{})**: A downward gravitational force acting at the center of mass of the rod, labeled "Gravity acts at the center of mass."
   - **\textbackslash{}( \textbackslash{}vec\{n\}\_1 \textbackslash{})**: A normal force acting upward at the base of the rod.
   - **\textbackslash{}( \textbackslash{}vec\{n\}\_2 \textbackslash{})**: A normal force acting horizontally where the rod contacts the wall.
   - **\textbackslash{}( \textbackslash{}vec\{f\}\_s \textbackslash{})**: A static frictional force acting horizontally to the left at the base of the rod, labeled "Static friction prevents slipping."

3. **Geometric Labels**:
   - **\textbackslash{}( d\_1 \textbackslash{})**: Distance from the base to where the static friction acts.
   - **\textbackslash{}( d\_2 \textbackslash{})**: Vertical distance from the ground to where the rod contacts the wall.
   - The center of mass is marked with a circle and cross symbol \textbackslash{}(\textbackslash{}otimes\textbackslash{}).

4. **Additional Annotations**:
   - Text indicating, "\textbackslash{}(\textbackslash{}tau\_\{\textbackslash{}text\{net\}\} = 0\textbackslash{}) about this point," suggesting torque is balanced around the point where the rod contacts the ground.
   - The angle \textbackslash{}( 60\textasciicircum{}\textbackslash{}circ \textbackslash{}) is shown between the rod and axial force lines at the base.

These elements illustrate the balance of forces and torques maintaining equilibrium for the rod in the system.

\#\# Dataset metadata

Statics; Physical Model Grounding Reasoning, Spatial Relation Reasoning

\paragraph{Recorded answer/context.}
C: C: 0.29

\paragraph{Figure/image path.}
/root/proposal\_for\_physic/science-mango/phyx\_mini\_run/phyx\_data/test\_image/709.png

\paragraph{Formalization target.}
create a compiling Lean file with sorry bodies at `PhyXMiniProblems/problem\_phyx\_mini\_0709.lean`. The Lean declarations must preserve the physical quantities, dimensions or dimensional roles, figure labels, governing-law hypotheses, and final relation expressed by this problem.
Use Mathlib/Physlib names found through LeanExplore where available. If a physics API is missing, introduce faithful local abstractions rather than scalar placeholder aliases.

\begin{theorem}[Physics formalization target]
\label{thm:physics:phyx_mini_0709:target}
\lean{PhyXMiniProblems.ProblemPhyXMini0709.minimumStaticFrictionCoefficient_is_answerChoiceC}
\uses{def:physics:phyx-mini-0709:phyxminiproblems-problemphyxmini0709-ladderstaticssetup, def:physics:phyx-mini-0709:phyxminiproblems-problemphyxmini0709-matchesproblemandprimaryfigure, def:physics:phyx-mini-0709:phyxminiproblems-problemphyxmini0709-hasuniformladdergeometry, def:physics:phyx-mini-0709:phyxminiproblems-problemphyxmini0709-hasphysicalladderparameters, def:physics:phyx-mini-0709:phyxminiproblems-problemphyxmini0709-satisfiesstaticforceandtorquebalance, def:physics:phyx-mini-0709:phyxminiproblems-problemphyxmini0709-coefficientpreventsslip, def:physics:phyx-mini-0709:phyxminiproblems-problemphyxmini0709-recordedanswerchoice, def:physics:phyx-mini-0709:phyxminiproblems-problemphyxmini0709-isclosestanswerchoice}
The assigned autoformalize agent should translate this physics problem into Lean declarations in the covered file, with theorem and lemma proof bodies written as `by sorry`.
\end{theorem}
\begin{proof}
This is an autoformalization task, not a proof task. Produce statements that can later be proved by the physics prover without weakening the physical model.
\end{proof}

% --- Archon declaration topology begin ---
\section{Declaration topology}
\begin{definition}[Length Magnitude]
  \label{def:physics:phyx-mini-0709:phyxminiproblems-problemphyxmini0709-lengthmagnitude}
  \lean{PhyXMiniProblems.ProblemPhyXMini0709.LengthMagnitude}
  A nonnegative physical length
\end{definition}
\begin{proof}
  This is the declared model definition of Length Magnitude.
\end{proof}
\begin{definition}[Force Magnitude]
  \label{def:physics:phyx-mini-0709:phyxminiproblems-problemphyxmini0709-forcemagnitude}
  \lean{PhyXMiniProblems.ProblemPhyXMini0709.ForceMagnitude}
  A nonnegative force magnitude, with dimension mass times acceleration
\end{definition}
\begin{proof}
  This is the declared model definition of Force Magnitude.
\end{proof}
\begin{definition}[length In Meters]
  \label{def:physics:phyx-mini-0709:phyxminiproblems-problemphyxmini0709-lengthinmeters}
  \lean{PhyXMiniProblems.ProblemPhyXMini0709.lengthInMeters}
  \uses{def:physics:phyx-mini-0709:phyxminiproblems-problemphyxmini0709-lengthmagnitude}
  Read a physical length in coherent SI metres
\end{definition}
\begin{proof}
  This is the declared model definition of length In Meters.
\end{proof}
\begin{definition}[force In Newtons]
  \label{def:physics:phyx-mini-0709:phyxminiproblems-problemphyxmini0709-forceinnewtons}
  \lean{PhyXMiniProblems.ProblemPhyXMini0709.forceInNewtons}
  \uses{def:physics:phyx-mini-0709:phyxminiproblems-problemphyxmini0709-forcemagnitude}
  Read a physical force magnitude in coherent SI newtons
\end{definition}
\begin{proof}
  This is the declared model definition of force In Newtons.
\end{proof}
\begin{definition}[angle In Radians From Degrees]
  \label{def:physics:phyx-mini-0709:phyxminiproblems-problemphyxmini0709-angleinradiansfromdegrees}
  \lean{PhyXMiniProblems.ProblemPhyXMini0709.angleInRadiansFromDegrees}
  Convert an angle readout in degrees to its dimensionless radian value
\end{definition}
\begin{proof}
  This is the declared model definition of angle In Radians From Degrees.
\end{proof}
\begin{definition}[Force Label]
  \label{def:physics:phyx-mini-0709:phyxminiproblems-problemphyxmini0709-forcelabel}
  \lean{PhyXMiniProblems.ProblemPhyXMini0709.ForceLabel}
  The four force-vector labels printed in the supplied figure
\end{definition}
\begin{proof}
  This is the declared model definition of Force Label.
\end{proof}
\begin{definition}[Force Direction]
  \label{def:physics:phyx-mini-0709:phyxminiproblems-problemphyxmini0709-forcedirection}
  \lean{PhyXMiniProblems.ProblemPhyXMini0709.ForceDirection}
  Cardinal directions of the force arrows in the supplied figure
\end{definition}
\begin{proof}
  This is the declared model definition of Force Direction.
\end{proof}
\begin{definition}[Surface Condition]
  \label{def:physics:phyx-mini-0709:phyxminiproblems-problemphyxmini0709-surfacecondition}
  \lean{PhyXMiniProblems.ProblemPhyXMini0709.SurfaceCondition}
  Contact idealizations for the two supporting surfaces
\end{definition}
\begin{proof}
  This is the declared model definition of Surface Condition.
\end{proof}
\begin{definition}[Ladder Point]
  \label{def:physics:phyx-mini-0709:phyxminiproblems-problemphyxmini0709-ladderpoint}
  \lean{PhyXMiniProblems.ProblemPhyXMini0709.LadderPoint}
  Named points on the ladder used by the force and torque annotations
\end{definition}
\begin{proof}
  This is the declared model definition of Ladder Point.
\end{proof}
\begin{definition}[Ladder Statics Figure]
  \label{def:physics:phyx-mini-0709:phyxminiproblems-problemphyxmini0709-ladderstaticsfigure}
  \lean{PhyXMiniProblems.ProblemPhyXMini0709.LadderStaticsFigure}
  \uses{def:physics:phyx-mini-0709:phyxminiproblems-problemphyxmini0709-forcelabel, def:physics:phyx-mini-0709:phyxminiproblems-problemphyxmini0709-forcedirection, def:physics:phyx-mini-0709:phyxminiproblems-problemphyxmini0709-ladderpoint}
  Literal labels and qualitative information visible in the primary raster. The figure does not print numerical values for d₁ or d₂; it only names them, so their geometry is imposed separately below
\end{definition}
\begin{proof}
  This is the declared model definition of Ladder Statics Figure.
\end{proof}
\begin{definition}[Ladder Statics Setup]
  \label{def:physics:phyx-mini-0709:phyxminiproblems-problemphyxmini0709-ladderstaticssetup}
  \lean{PhyXMiniProblems.ProblemPhyXMini0709.LadderStaticsSetup}
  \uses{def:physics:phyx-mini-0709:phyxminiproblems-problemphyxmini0709-lengthmagnitude, def:physics:phyx-mini-0709:phyxminiproblems-problemphyxmini0709-forcemagnitude, def:physics:phyx-mini-0709:phyxminiproblems-problemphyxmini0709-forcelabel, def:physics:phyx-mini-0709:phyxminiproblems-problemphyxmini0709-surfacecondition, def:physics:phyx-mini-0709:phyxminiproblems-problemphyxmini0709-ladderstaticsfigure}
  The dimensionful quantities of the statics problem. Each force field is a magnitude; its direction is recorded independently by the primary-figure data. In particular, no coefficient value or answer choice is stored here
\end{definition}
\begin{proof}
  This is the declared model definition of Ladder Statics Setup.
\end{proof}
\begin{definition}[Matches Problem And Primary Figure]
  \label{def:physics:phyx-mini-0709:phyxminiproblems-problemphyxmini0709-matchesproblemandprimaryfigure}
  \lean{PhyXMiniProblems.ProblemPhyXMini0709.MatchesProblemAndPrimaryFigure}
  \uses{def:physics:phyx-mini-0709:phyxminiproblems-problemphyxmini0709-lengthinmeters, def:physics:phyx-mini-0709:phyxminiproblems-problemphyxmini0709-angleinradiansfromdegrees, def:physics:phyx-mini-0709:phyxminiproblems-problemphyxmini0709-ladderstaticssetup}
  Problem and primary-image data: a 3.0 m ladder at 60°, the four force arrows with the displayed orientations, a frictionless wall, static friction at the ground, the two lever-arm labels, and the torque balance annotation at the ladder base. These data contain no proposed value of μ\_s
\end{definition}
\begin{proof}
  This is the declared model definition of Matches Problem And Primary Figure.
\end{proof}
\begin{definition}[Has Uniform Ladder Geometry]
  \label{def:physics:phyx-mini-0709:phyxminiproblems-problemphyxmini0709-hasuniformladdergeometry}
  \lean{PhyXMiniProblems.ProblemPhyXMini0709.HasUniformLadderGeometry}
  \uses{def:physics:phyx-mini-0709:phyxminiproblems-problemphyxmini0709-lengthinmeters, def:physics:phyx-mini-0709:phyxminiproblems-problemphyxmini0709-ladderstaticssetup}
  Geometry of the uniform ladder in the depicted branch. The center of mass is at the midpoint, so d₁ is the horizontal moment arm (L/2) cos θ; d₂ is the vertical wall-contact height L sin θ
\end{definition}
\begin{proof}
  This is the declared model definition of Has Uniform Ladder Geometry.
\end{proof}
\begin{definition}[Has Physical Ladder Parameters]
  \label{def:physics:phyx-mini-0709:phyxminiproblems-problemphyxmini0709-hasphysicalladderparameters}
  \lean{PhyXMiniProblems.ProblemPhyXMini0709.HasPhysicalLadderParameters}
  \uses{def:physics:phyx-mini-0709:phyxminiproblems-problemphyxmini0709-lengthinmeters, def:physics:phyx-mini-0709:phyxminiproblems-problemphyxmini0709-forceinnewtons, def:physics:phyx-mini-0709:phyxminiproblems-problemphyxmini0709-ladderstaticssetup}
  Positivity and angle conditions selecting the nondegenerate configuration shown in the image. They constrain only the physical branch, not the sought coefficient
\end{definition}
\begin{proof}
  This is the declared model definition of Has Physical Ladder Parameters.
\end{proof}
\begin{definition}[Satisfies Static Force And Torque Balance]
  \label{def:physics:phyx-mini-0709:phyxminiproblems-problemphyxmini0709-satisfiesstaticforceandtorquebalance}
  \lean{PhyXMiniProblems.ProblemPhyXMini0709.SatisfiesStaticForceAndTorqueBalance}
  \uses{def:physics:phyx-mini-0709:phyxminiproblems-problemphyxmini0709-lengthinmeters, def:physics:phyx-mini-0709:phyxminiproblems-problemphyxmini0709-forceinnewtons, def:physics:phyx-mini-0709:phyxminiproblems-problemphyxmini0709-ladderstaticssetup}
  Static translational and rotational equilibrium in the force directions shown by the figure. The horizontal balance is f\_s = n₂, the vertical balance is n₁ = F\_G, and zero torque about the base is n₂ d₂ = F\_G d₁
\end{definition}
\begin{proof}
  This is the declared model definition of Satisfies Static Force And Torque Balance.
\end{proof}
\begin{definition}[Coefficient Prevents Slip]
  \label{def:physics:phyx-mini-0709:phyxminiproblems-problemphyxmini0709-coefficientpreventsslip}
  \lean{PhyXMiniProblems.ProblemPhyXMini0709.CoefficientPreventsSlip}
  \uses{def:physics:phyx-mini-0709:phyxminiproblems-problemphyxmini0709-forceinnewtons, def:physics:phyx-mini-0709:phyxminiproblems-problemphyxmini0709-ladderstaticssetup}
  Coulomb's static-friction inequality for a proposed dimensionless coefficient μ. Such a coefficient can support the equilibrium when μ ≥ 0 and the required friction magnitude does not exceed μ n₁
\end{definition}
\begin{proof}
  This is the declared model definition of Coefficient Prevents Slip.
\end{proof}
\begin{definition}[Answer Choice]
  \label{def:physics:phyx-mini-0709:phyxminiproblems-problemphyxmini0709-answerchoice}
  \lean{PhyXMiniProblems.ProblemPhyXMini0709.AnswerChoice}
  Labels of the four answers supplied with the problem
\end{definition}
\begin{proof}
  This is the declared model definition of Answer Choice.
\end{proof}
\begin{definition}[coefficient]
  \label{def:physics:phyx-mini-0709:phyxminiproblems-problemphyxmini0709-answerchoice-coefficient}
  \lean{PhyXMiniProblems.ProblemPhyXMini0709.AnswerChoice.coefficient}
  \uses{def:physics:phyx-mini-0709:phyxminiproblems-problemphyxmini0709-answerchoice}
  Dimensionless static-friction coefficient printed beside each choice
\end{definition}
\begin{proof}
  This is the declared model definition of coefficient.
\end{proof}
\begin{definition}[recorded Answer Choice]
  \label{def:physics:phyx-mini-0709:phyxminiproblems-problemphyxmini0709-recordedanswerchoice}
  \lean{PhyXMiniProblems.ProblemPhyXMini0709.recordedAnswerChoice}
  \uses{def:physics:phyx-mini-0709:phyxminiproblems-problemphyxmini0709-answerchoice}
  The answer label recorded in the supplied dataset metadata
\end{definition}
\begin{proof}
  This is the declared model definition of recorded Answer Choice.
\end{proof}
\begin{definition}[Is Closest Answer Choice]
  \label{def:physics:phyx-mini-0709:phyxminiproblems-problemphyxmini0709-isclosestanswerchoice}
  \lean{PhyXMiniProblems.ProblemPhyXMini0709.IsClosestAnswerChoice}
  \uses{def:physics:phyx-mini-0709:phyxminiproblems-problemphyxmini0709-answerchoice}
  A displayed option is closest to an exact coefficient when its absolute error is no greater than that of any other displayed option
\end{definition}
\begin{proof}
  This is the declared model definition of Is Closest Answer Choice.
\end{proof}
% --- Archon declaration topology end ---
% --- Archon physics formalization source end ---
